\subsection{Architecture ablations}
\label{sec:ablation}

We evaluate the main architectural and implementation choices of DPA4 on the
same WBM-subsampled test set and report energy and force MAEs under a shared
evaluation protocol. Because the ablations isolate different factors only
within their own controlled groups, each group is interpreted as evidence for
the corresponding design choice rather than as part of a global ranking across
unrelated runs. Efficiency metrics are reported as relative quantities under
matched hardware, batch-size, data-loading, and precision settings. Unless
otherwise noted, \textit{Train time (rel.)} denotes relative training
wall-clock time, and \textit{Test time (rel.)} denotes the relative wall-clock
time required to evaluate the full WBM-subsampled test set with the DeePMD-kit
\texttt{dp test} command. Although the test set and hardware are matched, test
time is retained as a secondary throughput reference because model size changes
the feasible batch size, memory pressure, and kernel execution pattern.
Normalized training time is therefore the primary efficiency metric. In tables
where rank highlighting is used, boldface denotes the best value within the
corresponding controlled comparison and underlining denotes the second-best
value. We omit second-best underlining for metrics without an explicitly
highlighted best value or when multiple rows are tied for the best value.

\subsubsection{Compilation and training precision}

This ablation tests the effect of enabling \texttt{torch.compile} together with
the training-precision choices used by DPA4. We separate three factors: graph
compilation, bfloat16 automatic mixed precision (AMP), and TF32 tensor-core
matmul. As a graph-level execution transformation, \texttt{torch.compile} does
not change the model equations; correspondingly, the reported compiled
configurations do not show a visible accuracy loss attributable to compilation.
Table~\ref{tab:dpa4_compile_ablation} shows that bf16 AMP reduces the training
time by approximately 30\% and peak training memory by approximately 59\%
relative to the non-compiled FP32 baseline. In the compiled setting, bf16 AMP
preserves the WBM MAE of the FP32 run while reducing training time by
approximately 48\% and peak training memory by approximately 48\%. This
stability is consistent with the implementation: numerically sensitive
operations, including geometric preprocessing and normalization operations such
as RMSNorm, are kept in FP32 while large autocast-eligible matrix operations
use bf16. Combining \texttt{torch.compile} with bf16 AMP gives the largest
overall efficiency gain, with approximately 68\% lower training time and
approximately 60\% lower peak training memory than the baseline. Enabling TF32
on top of bf16 AMP provides no measurable efficiency benefit because most
matrix operations already execute in bf16 under AMP; the small MAE differences
between the TF32 and non-TF32 rows are within the expected run-to-run variation
and are not interpreted as a systematic accuracy effect.

\begin{table}[t]
    \centering
    \scriptsize
    \begin{threeparttable}
        \caption{\texttt{torch.compile} and training-precision ablations.\tnote{a}}
        \label{tab:dpa4_compile_ablation}
        \begin{tabular}{@{}cccccccc@{}}
            \toprule
            \textbf{Compile}                      &
            \textbf{AMP}                          &
            \textbf{TF32}                         &
            \textbf{E MAE}$\downarrow$            &
            \textbf{F MAE}$\downarrow$            &
            \textbf{Train (rel.)}$\downarrow$     &
            \textbf{Test time (rel.)}$\downarrow$ &
            \makecell{\textbf{Peak train}\\\textbf{mem. (rel.)}$\downarrow$}                                                                    \\
            \midrule
            False                                 & False & $\times$ & TBD             & TBD             & 1.00          & TBD  & 1.00          \\
            False                                 & True  & $\times$ & TBD             & TBD             & 0.70          & TBD  & 0.41          \\
            True                                  & False & False    & \textbf{28.012} & \textbf{37.458} & 0.62          & 1.02 & 0.77          \\
            True                                  & True  & False    & 28.108          & \underline{37.595} & \textbf{0.32} & 1.00 & \textbf{0.40} \\
            True                                  & True  & True     & \underline{28.051} & 38.224          & \textbf{0.32} & 1.01 & \textbf{0.40} \\
            \bottomrule
        \end{tabular}
        \begin{tablenotes}
            \item[a] Train time and peak train memory are normalized to the no-compile, no-AMP
            baseline measured on NVIDIA H20 hardware. Peak train memory denotes peak GPU
            memory during training. Test time is normalized to the compiled AMP/no-TF32
            row. The $\times$ symbol indicates that TF32 tensor cores are not applicable
            for the non-compiled execution path.
        \end{tablenotes}
    \end{threeparttable}
\end{table}

\subsubsection{Attention aggregation}

The attention ablation isolates the aggregation rule while holding the feature
dimension and focus count fixed within each pair of rows. Replacing Scatter Sum
with Attention Weighted Sum consistently reduces both energy and force MAEs
across the 64-channel, 96-channel, and 96-channel 2-focus settings. This
accuracy gain is obtained with only a modest computational overhead: training
time increases by approximately 5--6\% within each subgroup, and inference time
remains close to the corresponding Scatter Sum control. This result supports
the design choice of computing attention from rotationally invariant scalar
channels: the model gains adaptive neighbour selection while preserving
equivariance of the higher-order SO(2) features.

\begin{table}[t]
    \centering
    \scriptsize
    \begin{threeparttable}
        \caption{Attention aggregation ablation.\tnote{a}}
        \label{tab:dpa4_attention_ablation}
        \begin{tabular}{@{}ccccccc@{}}
            \toprule
            \textbf{Feature dim.}             &
            \textbf{No. focuses}              &
            \textbf{Aggregation}              &
            \textbf{E MAE}$\downarrow$        &
            \textbf{F MAE}$\downarrow$        &
            \textbf{Train (rel.)}$\downarrow$ &
            \textbf{Test time (rel.)}$\downarrow$                                                                            \\
            \midrule
            64                                & 1 & Scatter Sum            & \underline{30.691} & \underline{40.072} & 1.00 & 1.00 \\
            64                                & 1 & Attention Weighted Sum & \textbf{27.839} & \textbf{38.184} & 1.05 & 1.01 \\
            \midrule
            96                                & 1 & Scatter Sum            & \underline{30.068} & \underline{38.407} & 1.00 & 1.00 \\
            96                                & 1 & Attention Weighted Sum & \textbf{27.567} & \textbf{36.127} & 1.05 & 1.05 \\
            \midrule
            96                                & 2 & Scatter Sum            & \underline{30.935} & \underline{36.639} & 1.00 & 1.00 \\
            96                                & 2 & Attention Weighted Sum & \textbf{28.083} & \textbf{34.158} & 1.06 & 1.02 \\
            \bottomrule
        \end{tabular}
        \begin{tablenotes}
            \item[a] Within each feature-dimension and focus-count setting, train time and test
            time are normalized to the corresponding Scatter Sum aggregation.
        \end{tablenotes}
    \end{threeparttable}
\end{table}

\subsubsection{Multi-focus SO(2) convolution}

The multi-focus comparison separates per-focus feature width from the number of
parallel equivariant focus channels. Increasing the width of a single focus
stream rapidly increases parameter count and inference cost, but the
corresponding accuracy gains are uneven. By contrast, multi-focus variants
increase the effective SO(2) convolution dimension through several narrower
focus channels and often reach better accuracy--cost trade-offs. At an SO(2)
dimension of 192, the 96-channel 2-focus model gives the best overall result,
improving energy MAE from 29.418 to 26.994~meV/atom and force MAE from 39.529
to 36.408~meV/\AA{} relative to the 64-channel single focus baseline. It also
outperforms the 192-channel single focus model with more than 56\% fewer
trainable parameters, approximately 23\% lower training time, and approximately
34\% lower inference time. All rows in this sweep use the same learning-rate
setting; as larger-capacity variants often benefit from smaller tuned learning
rates, the reported MAEs of the largest configurations may be mildly
conservative. Under a shared training recipe, however, the 96-channel 2-focus
configuration gives the most balanced point in this sweep. These trends are
consistent with the intended role of focus competition, where parallel
equivariant sub-channels specialize to different edge-local geometric motifs
before the rotate-back step.

\begin{table}[t]
    \centering
    \scriptsize
    \begin{threeparttable}
        \caption{Feature-width and focus-count ablation for SO(2) convolution.\tnote{a}}
        \label{tab:dpa4_focus_ablation}
        \begin{tabular}{@{}cccccccc@{}}
            \toprule
            \textbf{Feature dim.}                 &
            \textbf{No. focuses}                  &
            \textbf{SO(2) dim.}                   &
            \textbf{E MAE}$\downarrow$            &
            \textbf{F MAE}$\downarrow$            &
            \textbf{Train (rel.)}$\downarrow$     &
            \textbf{Test time (rel.)}$\downarrow$ &
            \textbf{Params}                                                                                           \\
            \midrule
            64                                    & 1 & 64  & 29.418          & 39.529          & 1.00 & 1.00 & 1.9M  \\
            96                                    & 1 & 96  & 28.671          & 38.333          & 1.37 & 1.49 & 4.1M  \\
            64                                    & 2 & 128 & 28.126          & 38.123          & 1.55 & 1.55 & 3.2M  \\
            128                                   & 1 & 128 & 28.708          & 37.882          & 1.80 & 1.86 & 7.3M  \\
            64                                    & 3 & 192 & 27.890          & 37.283          & 1.95 & 2.14 & 4.5M  \\
            96                                    & 2 & 192 & \textbf{26.994} & \textbf{36.408} & 2.28 & 2.54 & 7.0M  \\
            192                                   & 1 & 192 & 27.286          & 36.477          & 2.98 & 3.86 & 16.0M \\
            64                                    & 4 & 256 & 27.821          & 37.605          & 2.43 & 2.65 & 5.7M  \\
            256                                   & 1 & 256 & 28.409          & 36.670          & 4.50 & 5.17 & 28.4M \\
            96                                    & 3 & 288 & \underline{27.168} & \underline{36.429} & 3.25 & 3.62 & 9.8M  \\
            \bottomrule
        \end{tabular}
        \begin{tablenotes}
            \item[a] Train time and test time are normalized to the 64-channel single focus
            baseline. Parameters are reported as trainable parameter counts. All rows use
            the same learning-rate setting.
        \end{tablenotes}
    \end{threeparttable}
\end{table}

\subsubsection{Edge-conditioned radial degree mixing}

Edge-conditioned radial degree mixing tests whether radial features should only
scale each angular degree independently or should also mix neighbouring degrees
in the local SO(2) frame. The diagonal modulation baseline multiplies each
coefficient by its degree-specific radial feature. The channel-shared dynamic
kernel instead mixes input degrees into output degrees with weights predicted
from the edge radial embedding and shared across channels. Sharing the kernel
across channels leaves the energy MAE essentially unchanged but reduces force
MAE from 38.349 to 37.212~meV/\AA{}. Introducing a low-rank channel-specific
kernel further improves the completed runs: rank 1 reduces force MAE to
35.689~meV/\AA{}, whereas rank 4 gives the best observed energy MAE of
26.556~meV/atom. The cost, however, increases with rank: rank 4 requires
1.86$\times$ the training time of the diagonal baseline. These results indicate
that dynamic radial degree mixing improves the force-sensitive angular
response, and that low-rank channel specificity provides additional accuracy at
a controllable but non-negligible training cost. Higher-rank and full kernels
require completed accuracy measurements before their accuracy--cost trade-offs
can be interpreted.

\begin{table}[t]
    \centering
    \scriptsize
    \begin{threeparttable}
        \caption{Ablation of edge-conditioned radial degree mixing.\tnote{a}}
        \label{tab:dpa4_radial_degree_ablation}
        \begin{tabular*}{\linewidth}{@{\extracolsep{\fill}}ccccccc@{}}
            \toprule
            \textbf{Radial Operation} &
            \textbf{Channel Specific} &
            \textbf{Rank} &
            \textbf{E MAE}$\downarrow$ &
            \textbf{F MAE}$\downarrow$ &
            \textbf{Train (rel.)}$\downarrow$ &
            \textbf{Test time (rel.)}$\downarrow$ \\
            \midrule
            Diagonal Modulation & None & -- & 28.493 & 38.349 & 1.00 & 1.00 \\
            Dynamic Degree Mixing & Shared & -- & 28.494 & 37.212 & 1.09 & 1.06 \\
            Dynamic Degree Mixing & Per Channel & 1 & 27.611 & \underline{35.689} & 1.12 & 1.10 \\
            Dynamic Degree Mixing & Per Channel & 2 & \underline{26.983} & 36.127 & 1.66 & 1.13 \\
            Dynamic Degree Mixing & Per Channel & 4 & \textbf{26.556} & \textbf{35.675} & 1.86 & 1.14 \\
            Dynamic Degree Mixing & Per Channel & 8 & TBD & TBD & 2.27 & 1.19 \\
            Dynamic Degree Mixing & Per Channel & 16 & TBD & TBD & 3.09 & TBD \\
            Dynamic Degree Mixing & Per Channel & Full & TBD & TBD & 3.36 & TBD \\
            \bottomrule
        \end{tabular*}
        \begin{tablenotes}
            \item[a] Train time and test time are normalized to the diagonal-modulation baseline.
            Rank 8 and larger Per Channel rows are timing-only placeholders and are not
            used for accuracy conclusions.
        \end{tablenotes}
    \end{threeparttable}
\end{table}

\subsubsection{S2 activation and quadrature}

This ablation separates two coupled design choices in the spherical-grid
nonlinearity: the model component to which S2 activation is applied and the
quadrature rule used to project grid features back to equivariant coefficients.
When S2 activation is restricted to the FFN component, replacing the e3nn
product grid with the Lebedev rule slightly reduces the energy MAE while leaving
the force MAE and computational cost essentially unchanged. Applying S2
activation additionally inside the SO(2) convolution component does not improve
the overall accuracy--cost trade-off under this setting; with the e3nn product
grid, both MAEs increase and the relative training and inference costs more than
double. In this expanded activation configuration, Lebedev quadrature reduces
the energy error and lowers the relative training and inference costs compared
with the e3nn product grid, although the force MAE remains higher than in the
FFN-only configuration. These results indicate that the main benefit comes from
retaining S2 activation in the FFN component, while Lebedev quadrature provides
a more favorable projection rule when spherical-grid nonlinearities are used in
equivariant branches.

\begin{table}[t]
    \centering
    \scriptsize
    \begin{threeparttable}
        \caption{Ablation of S2 activation placement and quadrature rule.\tnote{a}}
        \label{tab:dpa4_quadrature_ablation}
        \begin{tabular}{@{}ccccccc@{}}
            \toprule
            \textbf{SO(2) S2 Act.}            &
            \textbf{FFN S2 Act.}              &
            \textbf{Quadrature}               &
            \textbf{E MAE}$\downarrow$        &
            \textbf{F MAE}$\downarrow$        &
            \textbf{Train (rel.)}$\downarrow$ &
            \textbf{Test time (rel.)}$\downarrow$                                                                \\
            \midrule
            False                             & True & e3nn    & \underline{29.225} & \textbf{38.055} & 1.00 & 1.00 \\
            False                             & True & lebedev & \textbf{28.992} & \underline{38.103} & 0.98 & 1.01 \\
            True                              & True & e3nn    & 31.074          & 39.522          & 2.41 & 2.23 \\
            True                              & True & lebedev & 29.709          & 39.808          & 1.98 & 1.76 \\
            \bottomrule
        \end{tabular}
        \begin{tablenotes}
            \item[a] Train time and test time are normalized to the FFN-only S2 row using the
            e3nn product grid. The quadrature rule is used when projecting S2-grid features
            back to equivariant coefficients.
        \end{tablenotes}
    \end{threeparttable}
\end{table}

\subsubsection{Supplementary ablation studies}

Additional controlled studies are provided in the Supplementary Information to
document model selection and robustness analyses beyond the main mechanisms
tested above. These include model-depth scaling over the number of interaction
blocks and SO(2) layers per block, attention-aggregation design variants,
normalization-placement variants, and the learning-rate scheduler comparison.
These results support reproducibility and hyperparameter selection, whereas the
main text focuses on ablations that directly test the principal architectural
and implementation choices of DPA4.

